\clearpage
\begin{summarybox}
\textbf{\textcolor{CSummary}{一句话快记}}

\[
\boxed{\ f(x+a)=\frac{k}{f(x)}\ (k\ne0)\ }
\quad\Longrightarrow\quad
\boxed{\ f(x+2a)=f(x)\ }
\]
\textbf{记住：倒过来一次，再倒一次，就回来了。正周期写$2|a|$，不保证是最小正周期。}

\medskip
\textbf{\textcolor{CSummary}{三种做法，选一个就够}}
\begin{enumerate}[leftmargin=*,itemsep=0.35em]
\item \textbf{看数字（最好懂）：} $y\to k/y\to y$。
\item \textbf{连续代入（最常规）：} $f(x+2a)=k/f(x+a)=k/(k/f(x))=f(x)$。
\item \textbf{比较乘积（最快速）：} $f(x)f(x+a)=k=f(x+a)f(x+2a)$，消去中间因子。
\end{enumerate}

\medskip
\textbf{\textcolor{CSummary}{考场检查四句话}}
\begin{itemize}[leftmargin=*,itemsep=0.35em]
\item $k$是\textbf{同一个非零常数}，规则不能每次变来变去。
\item $f(x)$始终不能为$0$，所需位置必须在定义域内。
\item 先得到\textbf{一个周期}，没有证据不要写“最小正周期”。
\item 如果$a<0$，报告正周期时写$2|a|$。
\end{itemize}

\medskip
\textbf{\textcolor{CSummary}{最后两个值得记住的联系}}

\textbf{一是}“相邻两值的乘积不变”$f(x)f(x+a)=k$；\textbf{二是}负倒数型会使间隔$a$的函数值一正一负，因此不可能在整个实数范围上处处连续。这两点都能在高中的函数题中帮助你迅速看清题目。
\end{summarybox}
